\chapter{Kaula's inclination and eccentricity functions}\label{app:kaula}

\section{The transformation and its two functions}
Kaula's transformation \citep[\S3.3]{Kaula2000} converts a term of a spherical-harmonic potential of degree $l$ and order $m$, evaluated at the position of an orbiting body, into a double Fourier series in the body's angles. Two functions appear:
\begin{itemize}
  \item the inclination function $F_{lmp}(i)$, which results from rotating the frame from the equator to the orbital plane;
  \item the eccentricity function $G_{lpq}(e)$, which results from replacing $(r,f)$ by $(a,\M)$, i.e.\ the Hansen coefficients.
\end{itemize}
The real form of the transformation is \eqref{eq:kaula-real-CS}, the complex form \eqref{eq:kaula-complex}, and the spherical-harmonic form, with the factor $(-1)^m$ of the Condon--Shortley convention, is \eqref{eq:kaula-Ylm}. The angle is $\nu_{lmpq}=(l-2p)\omega+(l-2p+q)\M+m(\Omega-\thetap)$; the parity phase is $f_{lm}$.

\paragraph{Indices.} $l\ge2$ is the degree, $0\le m\le l$ the order, $0\le p\le l$ and $q\in\mathbb Z$ the indices of the inclination and eccentricity expansions. A second set $(k,j)$ is used when the transformation is applied at the field point (Chapter~\ref{ch:second}). Negative orders $m$ never appear in the results.

\section{Inclination function $F_{lmp}(i)$}
\paragraph{Definition} \citep[eq.~3.62]{Kaula2000}:
\begin{multline}
  F_{lmp}(i)=\sum_t\frac{(2l-2t)!}{t!(l-t)!(l-m-2t)!\,2^{2l-2t}}\sin^{l-m-2t}i\\
  \times\sum_{s=0}^m\binom ms\cos^si\sum_c\binom{l-m-2t+s}{c}\binom{m-s}{p-t-c}(-1)^{c-k}. \label{eq:F-def}
\end{multline}
Here $k=\lfloor(l-m)/2\rfloor$, $t$ runs from 0 to $\min(p,k)$, and $c$ over all values that make the binomials nonzero. The function is defined with Kaula's $\PK_{lm}$ (no Condon--Shortley phase).

\paragraph{Properties.}
\begin{itemize}
  \item Symmetry: $F_{l,m,l-p}(i)=(-1)^{l-m}F_{lmp}(\pi-i)$.
  \item Coplanar value: $F_{lmp}(0)=\PK_{lm}(0)$ if $l-m$ is even and $p=(l-m)/2$, and $0$ otherwise.
\end{itemize}
Tables~\ref{tab:F2mp} and \ref{tab:F3mp} give the exact closed forms, generated symbolically from \eqref{eq:F-def}. They agree with Kaula's Table~1.

% Generated from the exact series of the Kaula functions -- do not edit by hand
\begin{table}[ht]\centering\small
\caption{Kaula inclination functions $F_{2mp}(i)$ (Kaula 2000, eq.~3.62; exact). Rows with $F_{2mp}(0)\neq0$ are the only ones in the coplanar case.}
\label{tab:F2mp}
{\renewcommand{\arraystretch}{2.6}\begin{tabular}{ccc|>{$\displaystyle}l<{$}|>{$\displaystyle}l<{$}|>{$\displaystyle}c<{$}}\hline
$l$ & $m$ & $p$ & F_{lmp}(i) & \mathrm dF_{lmp}/\mathrm di & F_{lmp}(0)\\\hline
2 & 0 & 0 & - \frac{3 \sin^{2}{i}}{8} & - \frac{3 \sin{\left(2 i \right)}}{8} & 0\\
2 & 0 & 1 & \frac{3 \sin^{2}{i} - 2}{4} & \frac{3 \sin{\left(2 i \right)}}{4} & - \frac{1}{2}\\
2 & 0 & 2 & - \frac{3 \sin^{2}{i}}{8} & - \frac{3 \sin{\left(2 i \right)}}{8} & 0\\
\hline
2 & 1 & 0 & \frac{3 \left(\cos{i} + 1\right) \sin{i}}{4} & \frac{3 \left(\cos{i} + \cos{\left(2 i \right)}\right)}{4} & 0\\
2 & 1 & 1 & - \frac{3 \sin{\left(2 i \right)}}{4} & - \frac{3 \cos{\left(2 i \right)}}{2} & 0\\
2 & 1 & 2 & \frac{3 \left(\cos{i} - 1\right) \sin{i}}{4} & \frac{3 \left(- \cos{i} + \cos{\left(2 i \right)}\right)}{4} & 0\\
\hline
2 & 2 & 0 & \frac{3 \left(\cos{i} + 1\right)^{2}}{4} & - \frac{3 \left(\cos{i} + 1\right) \sin{i}}{2} & 3\\
2 & 2 & 1 & \frac{3 \sin^{2}{i}}{2} & \frac{3 \sin{\left(2 i \right)}}{2} & 0\\
2 & 2 & 2 & \frac{3 \left(\cos{i} - 1\right)^{2}}{4} & - \frac{3 \left(\cos{i} - 1\right) \sin{i}}{2} & 0\\
\hline
\end{tabular}}\end{table}

\begin{table}[ht]\centering\small
\caption{Kaula inclination functions $F_{3mp}(i)$ (exact; reference only, not needed by the single form).}
\label{tab:F3mp}
{\renewcommand{\arraystretch}{2.6}\begin{tabular}{ccc|>{$\displaystyle}l<{$}|>{$\displaystyle}c<{$}}\hline
$l$ & $m$ & $p$ & F_{lmp}(i) & F_{lmp}(0)\\\hline
3 & 0 & 0 & - \frac{5 \sin^{3}{i}}{16} & 0\\
3 & 0 & 1 & \frac{3 \left(5 \sin^{2}{i} - 4\right) \sin{i}}{16} & 0\\
3 & 0 & 2 & - \frac{3 \left(5 \sin^{2}{i} - 4\right) \sin{i}}{16} & 0\\
3 & 0 & 3 & \frac{5 \sin^{3}{i}}{16} & 0\\
\hline
3 & 1 & 0 & - \frac{15 \left(\cos{i} + 1\right) \sin^{2}{i}}{16} & 0\\
3 & 1 & 1 & - \frac{3 \left(\cos{i} + 10 \cos{\left(2 i \right)} + 15 \cos{\left(3 i \right)} + 6\right)}{64} & - \frac{3}{2}\\
3 & 1 & 2 & \frac{3 \left(20 \sin^{2}{i} + \cos{i} + 15 \cos{\left(3 i \right)} - 16\right)}{64} & 0\\
3 & 1 & 3 & \frac{15 \left(\cos{i} - 1\right) \sin^{2}{i}}{16} & 0\\
\hline
3 & 2 & 0 & \frac{15 \left(\cos{i} + 1\right)^{2} \sin{i}}{8} & 0\\
3 & 2 & 1 & - \frac{15 \left(\cos{i} + 1\right) \left(3 \cos{i} - 1\right) \sin{i}}{8} & 0\\
3 & 2 & 2 & \frac{15 \left(\cos{i} - 1\right) \left(3 \cos{i} + 1\right) \sin{i}}{8} & 0\\
3 & 2 & 3 & - \frac{15 \left(\cos{i} - 1\right)^{2} \sin{i}}{8} & 0\\
\hline
3 & 3 & 0 & \frac{15 \left(\cos{i} + 1\right)^{3}}{8} & 15\\
3 & 3 & 1 & \frac{45 \left(\cos{i} + 1\right) \sin^{2}{i}}{8} & 0\\
3 & 3 & 2 & - \frac{45 \left(\cos{i} - 1\right) \sin^{2}{i}}{8} & 0\\
3 & 3 & 3 & - \frac{15 \left(\cos{i} - 1\right)^{3}}{8} & 0\\
\hline
\end{tabular}}\end{table}


\section{Eccentricity function $G_{lpq}(e)$}
\paragraph{Definition.} $G_{lpq}$ is defined by Kaula's replacement \citep[p.~35]{Kaula2000}:
\begin{equation}
  \Big(\frac ar\Big)^{l+1}e^{\I(l-2p)f}=\sum_{q\in\mathbb Z}G_{lpq}(e)\,e^{\I(l-2p+q)\M}, \label{eq:G-def}
\end{equation}
i.e.\ $G_{lpq}$ is a Hansen coefficient:
\[
  G_{lpq}(e)=\frac1{2\pi}\int_0^{2\pi}\Big(\frac ar\Big)^{l+1}e^{\I(l-2p)f}e^{-\I(l-2p+q)\M}\,\dd\M.
\]

\paragraph{Properties.}
\begin{itemize}
  \item $G_{lpq}$ is real.
  \item Symmetry: $G_{l,l-p,-q}=G_{lpq}$.
  \item $G_{lpq}=\mathcal O(e^{|q|})$.
\end{itemize}

\paragraph{Tisserand's formula} \citep[eq.~3.67--3.69]{Kaula2000}, after \citet[p.~256]{Tisserand1889}. It holds for any $q$:
\begin{equation}
  G_{lpq}(e)=(-1)^{|q|}(1+\beta^2)^l\beta^{|q|}\sum_{k\ge0}P_{lpqk}Q_{lpqk}\beta^{2k},\qquad \beta=\frac{e}{1+\sqrt{1-e^2}}, \label{eq:tisserand}
\end{equation}
\[
  P_{lpqk}=\sum_{r=0}^{h}\binom{2p'-2l}{h-r}\frac{(-1)^r}{r!}\Big(\frac{(l-2p'+q')e}{2\beta}\Big)^r,\qquad h=\begin{cases}k+q',&q'>0,\\k,&q'\le0,\end{cases}
\]
\[
  Q_{lpqk}=\sum_{r=0}^{h}\binom{-2p'}{h-r}\frac{1}{r!}\Big(\frac{(l-2p'+q')e}{2\beta}\Big)^r,\qquad h=\begin{cases}k,&q'>0,\\k-q',&q'\le0,\end{cases}
\]
where:
\begin{itemize}
  \item $p'=p$, $q'=q$ for $p\le l/2$, and $p'=l-p$, $q'=-q$ for $p>l/2$;
  \item the binomials are generalised, $\binom nj=n(n-1)\cdots(n-j+1)/j!$ for any integer $n$.
\end{itemize}

\paragraph{Convergence.} The series converges geometrically in $\beta^2\approx e^2/4$. Number of $k$-terms for an absolute error below $10^{-13}$ on all $G_{2pq}$ needed by the truncation table of Chapter~\ref{ch:nbody}:
\begin{center}
\begin{tabular}{l|ccccccc}\toprule
$e$ & 0.01 & 0.05 & 0.1 & 0.2 & 0.3 & 0.4 & 0.5\\\midrule
$K$ & 4 & 5 & 6 & 8 & 11 & 14 & 22\\\bottomrule
\end{tabular}
\end{center}

\paragraph{Closed forms.} For $l-2p+q=0$ \citep[eq.~3.66]{Kaula2000}, for example $G_{210}=(1-e^2)^{-3/2}$ and $G_{31,-1}=e(1-e^2)^{-5/2}$. Also $G_{20,-2}=G_{222}=0$ exactly.

\paragraph{Power series.} Tables~\ref{tab:G20q}--\ref{tab:G21q} give the exact power series in $e$ for $l=2$, $|q|\le24$, obtained by expanding \eqref{eq:tisserand} in exact rational arithmetic. They extend Kaula's Table~2 and the tables of \citet{Cayley1861}. The series to $e^{40}$, with exact rational coefficients, are available as a data file accompanying this document.

Accuracy of the truncated series, as the maximum over the $G_{2pq}$ needed by the table:
\begin{center}
\begin{tabular}{l|cccccc}\toprule
order $N$ $\backslash$ $e$ & 0.05 & 0.1 & 0.2 & 0.3 & 0.4 & 0.5\\\midrule
7 & $3\times10^{-9}$ & $3\times10^{-6}$ & $8\times10^{-4}$ & $2\times10^{-2}$ & $0.2$ & $0.9$\\
15 & $3\times10^{-16}$ & $4\times10^{-15}$ & $9\times10^{-8}$ & $7\times10^{-5}$ & $6\times10^{-3}$ & $0.2$\\
20 & -- & -- & $8\times10^{-12}$ & $2\times10^{-6}$ & $6\times10^{-4}$ & $6\times10^{-2}$\\
30 & -- & -- & -- & $7\times10^{-12}$ & $6\times10^{-6}$ & $5\times10^{-3}$\\\bottomrule
\end{tabular}
\end{center}
A series truncated at $e^{30}$ is accurate to $7\times10^{-12}$ up to $e=0.3$ (for the reference row $10^{-5}$ of the truncation table of Chapter~\ref{ch:nbody}; the data file stops at $|q|=30$, and the series is taken as zero beyond). Beyond that, Tisserand's formula should be evaluated directly.

\paragraph{Exact values.} In Tables~\ref{tab:G20q} and~\ref{tab:G21q}, fractions longer than 12 digits are printed as decimals with 8 significant digits (relative rounding below $5\times10^{-8}$). The exact fractions are provided with this document:
\begin{itemize}
  \item \texttt{G\_exact\_fractions.pdf}: the same terms as Tables~\ref{tab:G20q}--\ref{tab:G21q}, each coefficient as an exact fraction and in double precision;
  \item \texttt{data/G\_exact\_l2.txt}: the complete series of $G_{20q}$ and $G_{21q}$, $|q|\le30$, to $e^{40}$, as exact fractions;
  \item \texttt{data/G\_series.json}: the same in machine-readable form, also for $l=3$.
\end{itemize}

% Generated from the exact series of the Kaula functions -- do not edit by hand
% Requires \usepackage{longtable,array}. Exact rational coefficients; full series to e^40 in G_series.json
{\renewcommand{\arraystretch}{2.2}\begin{longtable}{ccc|>{$\displaystyle}l<{$}}
\caption{Kaula eccentricity functions $G_{20q}(e)$, exact series to $e^{10}$, $|q|\le24$ (Tisserand, Kaula 2000 eq.~3.67--3.69); for $|q|>6$ the three leading terms. Long fractions are printed with 8 significant digits; the exact fractions of all these terms are given in the separate annex \texttt{annex\_G\_exact/G\_exact\_fractions.pdf}. $G_{22q}=G_{20,-q}$.}\label{tab:G20q}\\
\hline $l$ & $p$ & $q$ & G_{lpq}(e)\\\hline\endfirsthead
\hline $l$ & $p$ & $q$ & G_{lpq}(e)\\\hline\endhead
2 & 0 & -24 & 1.5875184\!\times\!10^{1}\,e^{24} - 5.4928137\!\times\!10^{1}\,e^{26} + 9.8184961\!\times\!10^{1}\,e^{28} + \mathcal O(e^{29})\\
2 & 0 & -23 & 1.1881177\!\times\!10^{1}\,e^{23} - 3.8242537\!\times\!10^{1}\,e^{25} + 6.4178774\!\times\!10^{1}\,e^{27} + \mathcal O(e^{28})\\
2 & 0 & -22 & 8.8967914\!\times\!10^{0}\,e^{22} - 2.6496966\!\times\!10^{1}\,e^{24} + 4.1623123\!\times\!10^{1}\,e^{26} + \mathcal O(e^{27})\\
2 & 0 & -21 & 6.6657926\!\times\!10^{0}\,e^{21} - 1.8255182\!\times\!10^{1}\,e^{23} + 2.6760737\!\times\!10^{1}\,e^{25} + \mathcal O(e^{26})\\
2 & 0 & -20 & 4.9971867\!\times\!10^{0}\,e^{20} - 1.2492967\!\times\!10^{1}\,e^{22} + 1.7039920\!\times\!10^{1}\,e^{24} + \mathcal O(e^{25})\\
2 & 0 & -19 & 3.7485643\!\times\!10^{0}\,e^{19} - 8.4811268\!\times\!10^{0}\,e^{21} + 1.0735007\!\times\!10^{1}\,e^{23} + \mathcal O(e^{24})\\
2 & 0 & -18 & 2.8137062\!\times\!10^{0}\,e^{18} - 5.7014573\!\times\!10^{0}\,e^{20} + 6.6844033\!\times\!10^{0}\,e^{22} + \mathcal O(e^{23})\\
2 & 0 & -17 & 2.1133575\!\times\!10^{0}\,e^{17} - 3.7864321\!\times\!10^{0}\,e^{19} + 4.1102717\!\times\!10^{0}\,e^{21} + \mathcal O(e^{22})\\
2 & 0 & -16 & 1.5883604\!\times\!10^{0}\,e^{16} - 2.4759735\!\times\!10^{0}\,e^{18} + 2.4947899\!\times\!10^{0}\,e^{20} + \mathcal O(e^{21})\\
2 & 0 & -15 & 1.1945401\!\times\!10^{0}\,e^{15} - 1.5864986\!\times\!10^{0}\,e^{17} + 1.4955082\!\times\!10^{0}\,e^{19} + \mathcal O(e^{20})\\
2 & 0 & -14 & \frac{157464}{175175}e^{14} - \frac{78732}{79625}e^{16} + 8.8765920\!\times\!10^{-1}\,e^{18} + \mathcal O(e^{19})\\
2 & 0 & -13 & 6.7675995\!\times\!10^{-1}\,e^{13} - 5.9216495\!\times\!10^{-1}\,e^{15} + 5.2499250\!\times\!10^{-1}\,e^{17} + \mathcal O(e^{18})\\
2 & 0 & -12 & 5.0968645\!\times\!10^{-1}\,e^{12} - 3.3325652\!\times\!10^{-1}\,e^{14} + 3.1330314\!\times\!10^{-1}\,e^{16} + \mathcal O(e^{17})\\
2 & 0 & -11 & 3.8386802\!\times\!10^{-1}\,e^{11} - 1.6794226\!\times\!10^{-1}\,e^{13} + 1.9262608\!\times\!10^{-1}\,e^{15} + \mathcal O(e^{16})\\
2 & 0 & -10 & \frac{4096}{14175}e^{10} - \frac{2048}{31185}e^{12} + \frac{58624}{467775}e^{14} + \mathcal O(e^{15})\\
2 & 0 & -9 & 2.1719477\!\times\!10^{-1}\,e^{9} - 5.4298693\!\times\!10^{-3}\,e^{11} + 8.8173673\!\times\!10^{-2}\,e^{13} + \mathcal O(e^{14})\\
2 & 0 & -8 & \frac{729}{4480}e^{8} + \frac{243}{8960}e^{10} + \frac{24057}{358400}e^{12} + \mathcal O(e^{13})\\
2 & 0 & -7 & \frac{15625}{129024}e^{7} + 4.1628641\!\times\!10^{-2}\,e^{9} + 5.3875451\!\times\!10^{-2}\,e^{11} + \mathcal O(e^{12})\\
2 & 0 & -6 & \frac{4}{45}e^{6} + \frac{2}{45}e^{8} + \frac{11}{252}e^{10} + \mathcal O(e^{11})\\
2 & 0 & -5 & \frac{81}{1280}e^{5} + \frac{81}{2048}e^{7} + \frac{19359}{573440}e^{9} + \mathcal O(e^{11})\\
2 & 0 & -4 & \frac{1}{24}e^{4} + \frac{7}{240}e^{6} + \frac{131}{5760}e^{8} + \frac{4303}{241920}e^{10} + \mathcal O(e^{11})\\
2 & 0 & -3 & \frac{1}{48}e^{3} + \frac{11}{768}e^{5} + \frac{313}{30720}e^{7} + \frac{3355}{442368}e^{9} + \mathcal O(e^{11})\\
2 & 0 & -2 & 0\\
2 & 0 & -1 & -\frac{1}{2}e + \frac{1}{16}e^{3} - \frac{5}{384}e^{5} - \frac{143}{18432}e^{7} - \frac{9097}{1474560}e^{9} + \mathcal O(e^{11})\\
2 & 0 & 0 & 1 - \frac{5}{2}e^{2} + \frac{13}{16}e^{4} - \frac{35}{288}e^{6} - \frac{5}{576}e^{8} - \frac{49}{3600}e^{10} + \mathcal O(e^{11})\\
2 & 0 & 1 & \frac{7}{2}e - \frac{123}{16}e^{3} + \frac{489}{128}e^{5} - \frac{1763}{2048}e^{7} + \frac{13527}{163840}e^{9} + \mathcal O(e^{11})\\
2 & 0 & 2 & \frac{17}{2}e^{2} - \frac{115}{6}e^{4} + \frac{601}{48}e^{6} - \frac{1423}{360}e^{8} + \frac{48619}{69120}e^{10} + \mathcal O(e^{11})\\
2 & 0 & 3 & \frac{845}{48}e^{3} - \frac{32525}{768}e^{5} + \frac{208225}{6144}e^{7} - 1.3840795\!\times\!10^{1}\,e^{9} + \mathcal O(e^{11})\\
2 & 0 & 4 & \frac{533}{16}e^{4} - \frac{13827}{160}e^{6} + \frac{104127}{1280}e^{8} - \frac{146017}{3584}e^{10} + \mathcal O(e^{11})\\
2 & 0 & 5 & \frac{228347}{3840}e^{5} - \frac{3071075}{18432}e^{7} + 1.7969544\!\times\!10^{2}\,e^{9} + \mathcal O(e^{11})\\
2 & 0 & 6 & \frac{73369}{720}e^{6} - \frac{775727}{2520}e^{8} + \frac{7527241}{20160}e^{10} + \mathcal O(e^{11})\\
2 & 0 & 7 & 1.6942345\!\times\!10^{2}\,e^{7} - 5.5010469\!\times\!10^{2}\,e^{9} + 7.4011937\!\times\!10^{2}\,e^{11} + \mathcal O(e^{12})\\
2 & 0 & 8 & \frac{634085}{2304}e^{8} - 9.5714070\!\times\!10^{2}\,e^{10} + 1.4131763\!\times\!10^{3}\,e^{12} + \mathcal O(e^{13})\\
2 & 0 & 9 & 4.3887349\!\times\!10^{2}\,e^{9} - 1.6292003\!\times\!10^{3}\,e^{11} + 2.6171604\!\times\!10^{3}\,e^{13} + \mathcal O(e^{14})\\
2 & 0 & 10 & 6.8943970\!\times\!10^{2}\,e^{10} - 2.7227149\!\times\!10^{3}\,e^{12} + 4.7253268\!\times\!10^{3}\,e^{14} + \mathcal O(e^{15})\\
2 & 0 & 11 & 1.0696726\!\times\!10^{3}\,e^{11} - 4.4798495\!\times\!10^{3}\,e^{13} + 8.3503215\!\times\!10^{3}\,e^{15} + \mathcal O(e^{16})\\
2 & 0 & 12 & 1.6423080\!\times\!10^{3}\,e^{12} - 7.2729041\!\times\!10^{3}\,e^{14} + 1.4486915\!\times\!10^{4}\,e^{16} + \mathcal O(e^{17})\\
2 & 0 & 13 & 2.4990369\!\times\!10^{3}\,e^{13} - 1.1670721\!\times\!10^{4}\,e^{15} + 2.4735012\!\times\!10^{4}\,e^{17} + \mathcal O(e^{18})\\
2 & 0 & 14 & 3.7734131\!\times\!10^{3}\,e^{14} - 1.8537629\!\times\!10^{4}\,e^{16} + 4.1646006\!\times\!10^{4}\,e^{18} + \mathcal O(e^{19})\\
2 & 0 & 15 & 5.6593667\!\times\!10^{3}\,e^{15} - 2.9180345\!\times\!10^{4}\,e^{17} + 6.9257766\!\times\!10^{4}\,e^{19} + \mathcal O(e^{20})\\
2 & 0 & 16 & 8.4377060\!\times\!10^{3}\,e^{16} - 4.5565401\!\times\!10^{4}\,e^{18} + 1.1391686\!\times\!10^{5}\,e^{20} + \mathcal O(e^{21})\\
2 & 0 & 17 & 1.2513987\!\times\!10^{4}\,e^{17} - 7.0639971\!\times\!10^{4}\,e^{19} + 1.8553645\!\times\!10^{5}\,e^{21} + \mathcal O(e^{22})\\
2 & 0 & 18 & 1.8472543\!\times\!10^{4}\,e^{18} - 1.0880399\!\times\!10^{5}\,e^{20} + 2.9951264\!\times\!10^{5}\,e^{22} + \mathcal O(e^{23})\\
2 & 0 & 19 & 2.7153443\!\times\!10^{4}\,e^{19} - 1.6660308\!\times\!10^{5}\,e^{21} + 4.7963238\!\times\!10^{5}\,e^{23} + \mathcal O(e^{24})\\
2 & 0 & 20 & 3.9761965\!\times\!10^{4}\,e^{20} - 2.5374321\!\times\!10^{5}\,e^{22} + 7.6247006\!\times\!10^{5}\,e^{24} + \mathcal O(e^{25})\\
2 & 0 & 21 & 5.8024101\!\times\!10^{4}\,e^{21} - 3.8457302\!\times\!10^{5}\,e^{23} + 1.2040123\!\times\!10^{6}\,e^{25} + \mathcal O(e^{26})\\
2 & 0 & 22 & 8.4407170\!\times\!10^{4}\,e^{22} - 5.8024478\!\times\!10^{5}\,e^{24} + 1.8896090\!\times\!10^{6}\,e^{26} + \mathcal O(e^{27})\\
2 & 0 & 23 & 1.2243244\!\times\!10^{5}\,e^{23} - 8.7185858\!\times\!10^{5}\,e^{25} + 2.9488753\!\times\!10^{6}\,e^{27} + \mathcal O(e^{28})\\
2 & 0 & 24 & 1.7711759\!\times\!10^{5}\,e^{24} - 1.3050284\!\times\!10^{6}\,e^{26} + 4.5779468\!\times\!10^{6}\,e^{28} + \mathcal O(e^{29})\\
\hline\end{longtable}}

{\renewcommand{\arraystretch}{2.2}\begin{longtable}{ccc|>{$\displaystyle}l<{$}}
\caption{Kaula eccentricity functions $G_{21q}(e)$, exact series to $e^{10}$, $|q|\le24$ (Tisserand, Kaula 2000 eq.~3.67--3.69); for $|q|>6$ the three leading terms. Long fractions are printed with 8 significant digits; the exact fractions of all these terms are given in the separate annex \texttt{annex\_G\_exact/G\_exact\_fractions.pdf}. $G_{22q}=G_{20,-q}$.}\label{tab:G21q}\\
\hline $l$ & $p$ & $q$ & G_{lpq}(e)\\\hline\endfirsthead
\hline $l$ & $p$ & $q$ & G_{lpq}(e)\\\hline\endhead
2 & 1 & -24 & 3.9497719\!\times\!10^{3}\,e^{24} - 1.9988026\!\times\!10^{4}\,e^{26} + 4.9741788\!\times\!10^{4}\,e^{28} + \mathcal O(e^{29})\\
2 & 1 & -23 & 2.8594427\!\times\!10^{3}\,e^{23} - 1.3744333\!\times\!10^{4}\,e^{25} + 3.2527219\!\times\!10^{4}\,e^{27} + \mathcal O(e^{28})\\
2 & 1 & -22 & 2.0688306\!\times\!10^{3}\,e^{22} - 9.4183356\!\times\!10^{3}\,e^{24} + 2.1144842\!\times\!10^{4}\,e^{26} + \mathcal O(e^{27})\\
2 & 1 & -21 & 1.4958257\!\times\!10^{3}\,e^{21} - 6.4291475\!\times\!10^{3}\,e^{23} + 1.3656270\!\times\!10^{4}\,e^{25} + \mathcal O(e^{26})\\
2 & 1 & -20 & 1.0807498\!\times\!10^{3}\,e^{20} - 4.3698127\!\times\!10^{3}\,e^{22} + 8.7564646\!\times\!10^{3}\,e^{24} + \mathcal O(e^{25})\\
2 & 1 & -19 & 7.8024284\!\times\!10^{2}\,e^{19} - 2.9557342\!\times\!10^{3}\,e^{21} + 5.5698964\!\times\!10^{3}\,e^{23} + \mathcal O(e^{24})\\
2 & 1 & -18 & 5.6281179\!\times\!10^{2}\,e^{18} - 1.9882666\!\times\!10^{3}\,e^{20} + 3.5114386\!\times\!10^{3}\,e^{22} + \mathcal O(e^{23})\\
2 & 1 & -17 & 4.0559164\!\times\!10^{2}\,e^{17} - 1.3290452\!\times\!10^{3}\,e^{19} + 2.1916798\!\times\!10^{3}\,e^{21} + \mathcal O(e^{22})\\
2 & 1 & -16 & 2.9198810\!\times\!10^{2}\,e^{16} - 8.8191037\!\times\!10^{2}\,e^{18} + 1.3526485\!\times\!10^{3}\,e^{20} + \mathcal O(e^{21})\\
2 & 1 & -15 & 2.0996289\!\times\!10^{2}\,e^{15} - 5.8019774\!\times\!10^{2}\,e^{17} + 8.2431081\!\times\!10^{2}\,e^{19} + \mathcal O(e^{20})\\
2 & 1 & -14 & 1.5078685\!\times\!10^{2}\,e^{14} - 3.7781499\!\times\!10^{2}\,e^{16} + 4.9522083\!\times\!10^{2}\,e^{18} + \mathcal O(e^{19})\\
2 & 1 & -13 & 1.0813325\!\times\!10^{2}\,e^{13} - 2.4298749\!\times\!10^{2}\,e^{15} + 2.9278831\!\times\!10^{2}\,e^{17} + \mathcal O(e^{18})\\
2 & 1 & -12 & 7.7418999\!\times\!10^{1}\,e^{12} - 1.5388210\!\times\!10^{2}\,e^{14} + 1.7006222\!\times\!10^{2}\,e^{16} + \mathcal O(e^{17})\\
2 & 1 & -11 & 5.5325854\!\times\!10^{1}\,e^{11} - 9.5550955\!\times\!10^{1}\,e^{13} + 9.6915004\!\times\!10^{1}\,e^{15} + \mathcal O(e^{16})\\
2 & 1 & -10 & 3.9452758\!\times\!10^{1}\,e^{10} - 5.7801470\!\times\!10^{1}\,e^{12} + 5.4188878\!\times\!10^{1}\,e^{14} + \mathcal O(e^{15})\\
2 & 1 & -9 & 2.8063401\!\times\!10^{1}\,e^{9} - 3.3715196\!\times\!10^{1}\,e^{11} + 2.9828532\!\times\!10^{1}\,e^{13} + \mathcal O(e^{14})\\
2 & 1 & -8 & \frac{178331}{8960}e^{8} - 1.8620290\!\times\!10^{1}\,e^{10} + 1.6340336\!\times\!10^{1}\,e^{12} + \mathcal O(e^{13})\\
2 & 1 & -7 & \frac{432091}{30720}e^{7} - 9.3806244\!\times\!10^{0}\,e^{9} + 9.1335431\!\times\!10^{0}\,e^{11} + \mathcal O(e^{12})\\
2 & 1 & -6 & \frac{3167}{320}e^{6} - \frac{2187}{560}e^{8} + \frac{195147}{35840}e^{10} + \mathcal O(e^{11})\\
2 & 1 & -5 & \frac{1773}{256}e^{5} - \frac{4987}{6144}e^{7} + \frac{418315}{114688}e^{9} + \mathcal O(e^{11})\\
2 & 1 & -4 & \frac{77}{16}e^{4} + \frac{129}{160}e^{6} + \frac{899}{320}e^{8} + \frac{107279}{40320}e^{10} + \mathcal O(e^{11})\\
2 & 1 & -3 & \frac{53}{16}e^{3} + \frac{393}{256}e^{5} + \frac{24753}{10240}e^{7} + \frac{211557}{81920}e^{9} + \mathcal O(e^{11})\\
2 & 1 & -2 & \frac{9}{4}e^{2} + \frac{7}{4}e^{4} + \frac{141}{64}e^{6} + \frac{197}{80}e^{8} + \frac{62401}{23040}e^{10} + \mathcal O(e^{11})\\
2 & 1 & -1 & \frac{3}{2}e + \frac{27}{16}e^{3} + \frac{261}{128}e^{5} + \frac{14309}{6144}e^{7} + \frac{423907}{163840}e^{9} + \mathcal O(e^{11})\\
2 & 1 & 0 & 1 + \frac{3}{2}e^{2} + \frac{15}{8}e^{4} + \frac{35}{16}e^{6} + \frac{315}{128}e^{8} + \frac{693}{256}e^{10} + \mathcal O(e^{11})\\
\hline\end{longtable}}

\noindent Symmetries: $G_{l,l-p,-q}(e)=G_{lpq}(e)$; for $l=2$, $G_{21q}=G_{21,-q}$ and $G_{22q}=G_{20,-q}$.
Leading order: $G_{2pq}=\mathcal O(e^{|q|})$. Closed forms (Kaula eq.~3.66, $l-2p+q=0$): $G_{210}=(1-e^2)^{-3/2}$, $G_{20,-2}=G_{222}=0$.


\section{Erratum of earlier tables}
The exact tables reveal the following errors in earlier working tables. They are listed so that they are not reused.
\begin{center}\small
\begin{tabular}{lll}\toprule
Function & written & exact\\\midrule
$G_{20,+2}$, coefficient of $e^4$ & $-115/16$ & $-115/6$ (as in Kaula, Table~2)\\
$G_{22,-2}=G_{20,+2}$ & $\tfrac{17}{2}e^2-\tfrac{115}{16}e^3$ & $\tfrac{17}{2}e^2-\tfrac{115}{6}e^4$\\
$G_{31,-6}$ & $\tfrac{48203}{9240}e^6$ & $\tfrac{48203}{9216}e^6$\\
$G_{310}$ & $1+\tfrac12e^2+\tfrac{239}{64}e^4-\tfrac{3323}{576}e^6$ & $1+2e^2+\tfrac{239}{64}e^4+\tfrac{3323}{576}e^6$\\
$G_{31,q}$, $q=3..7$ & copied from $G_{30,q}$ & e.g.\ $G_{313}=\tfrac{77}{6}e^3-\tfrac{25}{48}e^5+\dots$\\
$F_{311}$ & $\tfrac{15}{16}\sin^2i(1+3\cos i)-\tfrac34(1+3\cos i)$ & $\tfrac{15}{16}\sin^2i(1+3\cos i)-\tfrac34(1+\cos i)$\\
$F_{312}$ & $\tfrac{15}{16}\sin^2i(1-3\cos i)-\tfrac34(1-3\cos i)$ & $\tfrac{15}{16}\sin^2i(1-3\cos i)-\tfrac34(1-\cos i)$\\
$F_{331}$, $F_{332}$ & $\tfrac{48}{8}\sin^2i(1\pm\cos i)$ & $\tfrac{45}{8}\sin^2i(1\pm\cos i)$\\
$\dd F_{201}/\dd i$ & $\tfrac32\sin i\cos i-\tfrac12$ & $\tfrac32\sin i\cos i$\\\bottomrule
\end{tabular}
\end{center}
